\documentclass{article}
\usepackage[utf8]{inputenc}
\usepackage{amsmath}
\usepackage{graphicx}
\title{A Sample Paper}
\author{Jane Doe \and John Smith}
\date{2026-03-23}
\begin{document}
\maketitle

\begin{abstract}
This paper demonstrates the LaTeX reader.
\end{abstract}

\section{Introduction}

This is the introduction with inline math $E = mc^2$ and a citation \cite{smith2020}.

\subsection{Background}

See Figure~\ref{fig:example} and Table~\ref{tab:results}.

\begin{figure}[htbp]
\centering
\includegraphics{example.png}
\caption{An example figure}
\label{fig:example}
\end{figure}

\begin{table}[htbp]
\centering
\begin{tabular}{|l|r|}
\hline
Name & Value \\
\hline
Pi & 3.14 \\
E & 2.72 \\
\hline
\end{tabular}
\caption{Results}
\label{tab:results}
\end{table}

\section{Methods}

We used \textbf{bold} and \emph{italic} formatting. The equation:

\begin{equation}
\nabla \cdot \mathbf{E} = \frac{\rho}{\epsilon_0}
\end{equation}

\begin{itemize}
\item First item
\item Second item
\end{itemize}

\begin{enumerate}
\item Step one
\item Step two
\end{enumerate}

\begin{quote}
A famous quote goes here.
\end{quote}

\begin{verbatim}
fn main() {
    println!("Hello!");
}
\end{verbatim}

\end{document}
